\begin{table}[htbp]\centering\footnotesize
\caption{Anexos transversales: totales y variación, línea P (mdp)}\label{cua:trans}
\begin{tabular}{lrrr}
\toprule
Anexo & 2026 & 2027 & Var. real (\%) \\
\midrule
13. Erogaciones para la igualdad entre mujeres y… & 599.145,4 & 636.133,1 & +2,1 \\
14. Recursos para la atención de grupos vulnerab… & 687.765,0 & 791.471,1 & +10,7 \\
15. Estrategia nacional de transición energética & 17.868,4 & 20.289,1 & +9,2 \\
16. Recursos para la adaptación y mitigación de … & 212.569,7 & 160.120,9 & $-$27,6 \\
17. Erogaciones para el desarrollo de los jóvenes & 735.177,8 & 730.916,0 & $-$4,4 \\
18. Recursos para la atención de niñas, niños y … & 1.099.514,0 & 1.101.551,7 & $-$3,7 \\
19. Acciones para la prevención del delito, comb… & 465.729,8 & 465.930,1 & $-$3,8 \\
31. Consolidación de una sociedad de cuidados & 466.674,9 & 475.959,7 & $-$1,9 \\
33. Diversidad sexual y de género & \textbf{no existía} & 43.369,4 & --- \\
\emph{Suma aritmética de los nueve} & --- & \emph{4.425.741,1} & --- \\
\bottomrule
\end{tabular}
\\[2pt]\begin{minipage}{0.92\textwidth}\scriptsize \textbf{La última fila no es un agregado: es una advertencia.} Los anexos etiquetan los mismos programas más de una vez, de modo que su suma ---4.425.741,1 mdp, el 61~\% del gasto programable pagado--- \textbf{no mide ningún gasto}. Fuente: anexos 13 a 19, 31 y 33 de los dos proyectos de decreto.\end{minipage}
\end{table}
